\Needspace{10\baselineskip}
\subsection{A.159. Standard bounds and higher PL comparison over an arbitrary real closed field}

\begin{saresult}{A.159}{Proved under stated hypotheses}{Higher surjectivity over every fixed o-minimal real closed field}
Let $M$ be any fixed o-minimal expansion of a real closed field,
let $B$ be a closed compact PL manifold over $M$, and put
$N=B\times I$. Then
\[
 \pi_k P_{\mathrm{PL}}(B)\longrightarrow
 \pi_k P_{\mathrm{def}}(B)
 \quad\hbox{is surjective for every }k\geq1.
 \tag{159.1}
\]
Together with A.113 this gives the component bijection and
surjectivity in every positive degree. This statement does not
assert injectivity in positive degrees over an arbitrary
nonreal model.

The sufficient regularity criteria of A.147 and A.148 are valid
over $M$ when applied on a fixed standard PL model and the
forward parameter Lipschitz bounds are bounded by ordinary
finite integers. In the pair criterion, the factors preserve
both spatial ends setwise and agree on the bottom and on the
parameter boundary. Their quotient is then definably homotopic,
through based bottom-fixed concordance families, to a jointly
PL family. No metric bound on either inverse is required.
\end{saresult}

\begin{proof}
Throughout this proof a \emph{standard bound} means a bound by
an ordinary integer, viewed in $M$. It does not mean merely
the existence of an element of $M$ bounding a function.
All simplicial complexes, posets, lists, and inductions below
are finite in the ordinary sense.
All coordinate and Hausdorff estimates use the maximum norm.

\textbf{The standard spatial model.}\quad
Choose a finite triangulation of $B$ and replace it by the
standard realization of its abstract simplicial complex,
with the vertices taken among the standard basis vectors.
The vertex correspondence is a PL homeomorphism over $M$,
so conjugation by its product with $\mathrm{id}_I$ identifies
the PL and definable concordance problems. All subsequent
metric estimates are taken on this standard model.

The real-algebraic realization is a closed PL manifold as well.
Indeed, the finite PL sphere or ball charts for the links in
the given realization have finite triangulations and affine
descriptors. Fixing those descriptors makes their existence
a first-order statement in real closed fields; the statement
holds over $M$, hence over the real algebraic numbers. This
uses finite PL chart certificates, not a comparison of
parameterized PL homotopies. Fix such charts once and for all.
The product $N$, its double $X=B\times S^1$, the two halves,
and the seams $A_0,A_1$ therefore have fixed standard finite
PL models. Any finite Lipschitz constants of maps with these
fixed real-algebraic PL descriptors have standard bounds.

We also use the cubical PL model of A.154. Explicitly, the
barycentre of a face $\sigma$ in the barycentric subdivision
is sent to its characteristic vector $\mathbf1_\sigma$.
The resulting affine maps on chains triangulate
\[
 Q=\{x\in[0,1]^v:
          \operatorname{supp}(x)\in K,\ \max_i x_i=1\}.
\]
Thus the conjugating map is PL, has standard Lipschitz bounds
in both directions, and takes every marked subcomplex to a
union of coordinate faces.

\textbf{Definable contractions require no completeness theorem.}\quad
Let $T:[0,1]^a\to[0,1]^a$ be continuous and definable, with
\[
 \|T(t)-T(t')\|_\infty\le c\|t-t'\|_\infty,\qquad c<1.
\]
The definable continuous residual
$r(t)=\|T(t)-t\|_\infty$ attains its minimum $m$ by definable
compactness. If $m>0$ and $t_0$ realizes it, then
\[
 r(T(t_0))
 =\|T(T(t_0))-T(t_0)\|_\infty
 \le c\,r(t_0)<m,
\]
a contradiction. Hence $T$ has a fixed point; contraction
makes it unique. For fixed points of $T$ and $T'$,
\[
 \|t-t'\|_\infty
 \le (1-c)^{-1}\sup_z\|T(z)-T'(z)\|_\infty.
 \tag{159.2}
\]
Consequently the fixed point is definable and varies
continuously in any jointly continuous definable compact
family with a common contraction constant. Definable uniform
continuity supplies the required local uniform convergence
in (159.2). No infinite iteration or ordinary Banach
completeness is used.

\textbf{Finite target choices with standard bounds.}\quad
Consider first one forward parameter-Lipschitz family $V$
on the standard model, with bound $L\in\mathbb N$.
The auxiliary choices in A.147 can be made independently of
$V$ and entirely on the fixed standard target geometry.
A finite marked triangulation and its barycentric coordinate
functions give the Hudson sequence
\[
 0=f^0\le f^1\le\cdots\le f^{R-1}=1.
\]
Consecutive differences are supported in prescribed vertex
stars. Subdividing the unit rise into an ordinary finite
number of increments makes the common spatial Lipschitz
constant $c$ of the functions satisfy $cL<1$.
The needed number of increments is an ordinary integer
because the original geometric constants and $L$ have
standard bounds. This is the step that cannot be justified
for an unrestricted infinite element $L\in M$.

All subsequent finite target data in A.147 depend only on
this sequence, the standard spatial model, and the parameter
dimension: the interval boxes, support-label poset, small
conic ball configurations, strict enlargements, marked
coarsenings, and the target coloring of Mn\"ev's Lemma~8
with its unprescribed parameter part. They can be fixed over
the real algebraic numbers and then interpreted in $M$.
One precise way to make this assertion is to perform the
finite pure PL construction over the reals, retain its
finite combinatorial types and affine descriptors, and
express all incidences, strict inclusions, ball charts, and
compatibilities by their finite semialgebraic certificates.
Real-closed-field completeness then supplies the same
certificates over the real algebraic numbers. This transfers
only auxiliary finite PL data, never the given definable
family or a purported PL homotopy of it.

The marked finite construction and its use are detailed next,
so that no non-Archimedean Lebesgue-number argument is hidden
in this choice of data.


\textbf{The graph deformation and finite support labels.}\quad
Extend $V$ to the marked double by reflection. The fixed PL
reflection changes its parameter bound only by a standard
factor, which we absorb into the integer $L$. Use the same
Hudson functions and the following formulas from A.147 and
A.156:
\[
 \rho_{y,v}(q)=\min_{0\leq j<R}
       \max\{f^j(y),q-2jv\},\qquad
 T_v=1+2(R-1)v,\quad D=2R-1,
\]
\[
 \Lambda(v,p)=\max\{p,D(p-1)+T_v\}.
\]
For $(x,p,v)$, solve the finite fixed-point equations
\[
 t_i=\rho_{V_t(x),v}(\Lambda(v,p_i)),\qquad
 H_{p,v}(x)=V_t(x).
 \tag{159.3}
\]
The right side is a contraction in $t$ with constant $cL<1$;
the preceding residual-minimum argument proves existence,
uniqueness, definability, and joint continuity over $M$.
The inverse formula is
\[
 H_{p,v}^{-1}(y)=
 V^{-1}_{(\rho_{y,v}(\Lambda(v,p_i)))_i}(y).
 \tag{159.4}
\]
It verifies the homeomorphism assertion directly. The initial
map is $V$, and all spatial marks are preserved. If $V$ fixes
the bottom or is the identity on a parameter face, so does
this deformation. These conclusions use the endpoint values
of $\rho$, not a compactness assertion about ordinary topology.

At $v=1$, put $q_i=Dp_i$. The moving bands are
$[2j,2j+1]$, $0\le j<R-1$. On such a band,
\[
 \rho_{y,1}(q_i)=
 \min\{f^{j+1}(y),\max\{f^j(y),q_i-2j\}\},
\]
and on the intervening plateau it equals $f^{j+1}(y)$;
the last index contributes only the final constant plateau.
Thus varying $q_i$ within that band can change (159.4) only
where $f^{j+1}-f^j$ is nonzero, in its prescribed vertex star.
Choose the finite standard interval cover used in A.147 so
that each interval meets at most one moving band. For a
coordinate box or a nonempty intersection of these boxes,
its active support is the union of at most $k$ such stars.
Define its label by the bands actually met by that intersection,
not by taking the union of the original box labels. Consequently
inclusion of boxes or their intersections gives inclusion of
support labels in the required direction. This is a finite
poset, and every geometric choice in it is independent of $V$.

The small-ball separation construction of Mn\"ev,
Proposition~57, pp.~89--90, as used in A.147, replaces those
monotone star supports by compatible configurations of finitely
many conic PL balls. Retain the full finite indexing poset,
even when two old configuration labels coincide. The explicit
finite rank induction of A.147 enlarges the balls to make
proper indexing relations strict. This is the finite conic
extension of the construction behind Proposition~35,
pp.~61--62; it is not an invocation of that proposition with
an unmet strictness hypothesis. All input complexes and the
finite support poset have already been fixed over the real
algebraic numbers. The resulting choices can therefore be
included in the finite target certificates described above.
No subdivision is being asked to beat an infinitesimal
Lebesgue number.

\textbf{The marked coarsening and the coloring.}\quad
All target complexes are chosen compatible with the two
halves and the seams $A_0,A_1$, and the common fine complex
refines their barycentric subdivisions. Proposition~36,
pp.~62--63, supplies the conic-star invariance and its
cell-incidence property. Since the seams are locally flat
subcomplexes, these properties imply that each consistent
ball is either disjoint from the seams or forms with its
intersection a standard ball/equatorial-disk pair. This is
an inference using local flatness, not an additional literal
assertion in Proposition~36. The incidence property requires
the center to lie in any seam met by the ball, so it cannot
meet both disjoint seams.
For a crossing ball replace the single interior cell of the
unmarked coarsening by its two half-ball interiors and the
middle-disk interior, retaining their boundary cells. This
is the split marked coarsening of A.147, rather than the
original inscribed-ball conclusion of Lemma~6.

The terminal inverse images of the full balls are constant
over the relevant parameter cells. Outside the balls the
inverse map is pointwise constant; continuity gives the same
assertion on their boundaries. Hence any chosen boundary
subcells also have constant inverse images. For a fixed
mark $A$ or half $N^{\pm}$, preservation gives the exact
identities
\[
 H_p^{-1}(D\cap A)=H_p^{-1}(D)\cap A,\qquad
 H_p^{-1}(D\cap N^{\pm})=H_p^{-1}(D)\cap N^{\pm}.
 \tag{159.5}
\]
Thus split prismaticity follows directly. In the decreasing
rank induction, boundary triangulations are chosen compatible
with the equatorial disks and extended on the complement by
finite relative PL triangulation. For nested balls, intersecting
with the same fixed global marks preserves every required
incidence. This is exactly the finite marked refinement checked
in A.147; it involves finite PL subdivisions and no limiting
process.

Apply the proof of Mn\"ev's Lemma~8, pp.~67--68, with its
unprescribed parameter part, no additional maps, and the
fixed target coloring. Its choices depend only on that coloring
and the configuration presheaf. They therefore produce the
same finite terminal target coloring for any families using
the same Hudson data. The equations that verify prismaticity
are the support constancies just proved, so their validity
over $M$ follows from the explicit formulas. The finite target
coloring itself may be taken from the real-algebraic auxiliary
certificates. This is the only use of the real version of the
finite PL construction; it supplies fixed finite data, not a
comparison theorem for the input family.

\textbf{Finite prism straightening over the field.}\quad
The target prisms are the literal products $C\times\sigma$
of Mn\"ev's Proposition~4, p.~21. Equation~(11) and
Proposition~5, pp.~22--23, identify their inverse images as
$D\times\sigma$, with $D$ independent of the parameter on
the entire closed simplex. These are finite ball complexes;
the whole boundary of each product ball is a union of proper
prisms of smaller total dimension. A lower parameter face
may carry a spatial refinement, which does not affect this
boundary assertion.

Triangulate the finitely many definable spatial source cells
compatibly, together with every cell of a fixed target PL
complex containing all marks. The finite marked
Hauptvermutung, Shiota, Theorem~2.1, p.~5, applies over the
given real closed field. If $\tau$ is the compatible definable
triangulation and $a$ is the marked PL homeomorphism supplied
by that theorem, then $\phi=\tau a^{-1}$ preserves every
cell of the fixed target complex and makes all source cells
PL after taking inverse images. Each is a PL ball by the
same theorem applied to its pair with its boundary.
The cell-preserving map $\phi$ is definably isotopic to the
identity by the finite cellwise Alexander construction below.
Thus common conjugation by $\phi$ does not change the
concordance class or any equality on the marked subcomplexes.

The PL endpoint is constructed by induction on total prism
dimension. Prescribe the required map on the already assigned
subcomplex. For a remaining product $D\times\sigma$, all of
$(\partial D\times\sigma)\cup(D\times\partial\sigma)$ is
already assigned. Choose compatible finite boundary
subdivisions making this map simplicial, without changing
any previously assigned map. Take fixed PL charts of $D$
and its target cell to convex balls and cone the TOTAL
product boundary from centers having the same base coordinate
$b_\sigma$. The resulting map is PL and preserves projection:
source and target points on corresponding cone segments have
the same parameter $(1-u)b_\sigma+up$. This is the construction
in Mn\"ev's Lemma~3, proof Steps~2--3, pp.~24--25.
It is a finite affine construction over $M$.

For completeness, the homotopy from a definable cellular map
to its PL endpoint also uses only finite field operations and
definable compactness. After taking the quotient by the PL
endpoint, construct the homotopy to the identity in increasing
prism dimension. In one spatial ball, let the parameter ball
be $D_0=\sigma\times I$. The already prescribed spatial
boundary maps extend radially to a continuous definable
family $R_p$ of homeomorphisms of the whole spatial ball.
On $p\in\partial D_0$, divide the prescribed whole-ball map
by $R_p$ to obtain a family $Q_p$ fixing the spatial boundary.
In fixed cube coordinates write $p=ru$, with
$r=\|p\|_\infty$ and $u\in\partial D_0$. Fill this family
using
\[
 A_r(Q_u)(x)=
 \begin{cases}
 rQ_u(x/r),&\|x\|_\infty\le r,\ r>0,\\
 x,&\|x\|_\infty\ge r,
 \end{cases}
 \qquad A_0=\mathrm{id}.
 \tag{159.6}
\]
The formulas agree on the overlap. Their displacement is at
most $2r$, proving joint continuity at the parameter center
in the field order topology. Multiplying by $R_p$ extends
all the prescribed data. Fixed PL changes of ball coordinates
handle general cells. Finite pasting gives the required joint
definable homotopy, relative to any already prescribed
subcomplex. This also proves the cellwise contraction used
for $\phi$. No assertion that this particular homotopy is
jointly PL is needed.

For a single based bottom-fixed input, prescribe the identity
on the bottom and on the parameter boundary before the prism
induction. This proves the standard-bound sufficient criterion
corresponding to A.147.

\textbf{The common-data pair criterion.}\quad
Now let $r,s$ be a pair as in the statement, with a common
standard forward bound. Choose the same Hudson sequence,
interval boxes, configurations, and target coloring for both.
In (159.3), the two fixed-point systems coincide for inputs
on the bottom, since the factors agree there. On a parameter
face, the corresponding coordinate $t_i$ is the fixed endpoint;
the remaining systems then evaluate only that face, on which
the factors agree. Hence both deformations agree throughout on
\[
 E=(A_0\times P)\cup(N\times\partial P).
 \tag{159.7}
\]
After reflection use the corresponding full doubled parameter
boundary. The configuration presheaf is literally the same
for both maps, so the unprescribed version of Lemma~8 gives
one terminal target coloring, not two unrelated choices.
Their source prism cells coincide on $E$ because the maps
coincide there. Choose one $\phi$ triangulating all source
cells of both families.

First construct one common PL endpoint on the subcomplex
$E$, and extend it over the two full prism complexes
separately by total product coning. The resulting PL
endpoints have the same restriction to $E$. Their quotients
with the two definable terminal maps consequently also
agree on $E$. Construct one common homotopy on $E$ first;
then extend that entire homotopy over each full prism complex
by (159.6). Thus equality on $E$ holds at every homotopy time.
Taking the quotient of the two families produces a based,
bottom-fixed concordance homotopy to the quotient of the
PL endpoints. This proves the sufficient pair criterion
corresponding to A.148 over $M$. Every source-dependent step
here is finite compatible triangulation, marked
Hauptvermutung, or explicit coning; none uses a mesh estimate.

Only the standard numerical bound was used to choose the
finite Hudson target data. In particular this argument has
not replaced a potentially infinite parameter bound by an
integer, and it has not inferred a parameter-preserving PL
homotopy from a path in a coefficient space.


\textbf{Producing the required standard bounds.}\quad
Let $U$ be an arbitrary based definable $k$-cube of concordances
on the standard model. The spatial factorization of A.154
works over $M$ and gives
\[
 U_p=s_pr_p^{-1}.
 \tag{159.8}
\]
Both factors preserve the spatial marks, agree on the bottom,
and are the identity on the parameter boundary. They have a
common STANDARD spatial Lipschitz bound. To check the last
word, run the coordinate clocks on the unit cubical model.
Every evolving component remains in a fixed unit cube. Each
stage has an ordinary finite monotonicity count $N$, so its
variation bound and all mixed-coordinate amplification factors
are ordinary finite numerical expressions in $N$, the number
of coordinates, and the preceding bounds. The final bounds
are therefore below an ordinary integer. The fixed PL changes
of spatial model multiply them only by standard constants.

Form the disjoint tagged graph family
\[
 \mathcal G_p=
  (\{0\}\times\Gamma r_p)\cup
  (\{3\}\times\Gamma s_p)
 \tag{159.9}
\]
on the unit cubical model. The tag separation exceeds the
diameter of either untagged graph space. Hence the Hausdorff
distance of these unions is the maximum of the two individual
graph distances. Divide all ambient coordinates by $3$ to put
the family in a unit cube; undoing this scaling costs only the
standard factor $3$. It is a continuous compact-set family and
is constant on the parameter boundary.

Apply A.158 to obtain a compact resolution consisting of
finitely many integer-tagged unit cubes, with standard
projection and Hausdorff constants. Every coordinate of this
compact source is bounded by an ordinary integer. Append the
boundary copy used in A.157 with an integer tag larger than
all existing tags and diameters. Because the boundary family
is constant, this preserves standard bounds for both the
projection and the pulled family.

The labelled join of A.157 uses at most $k+2$ slots. If its
source coordinate bound is $M_0$, its near-pair estimate is
$m(1+M_0)$ times the original Lipschitz bound, and its far-pair
estimate uses $2m$ times the bounded output diameter. Here
$m,M_0$, and all those input bounds are standard. Thus the
projection and compact-set family on the join have standard
Lipschitz constants. The finite universal resolution and
relative homotopy-extension argument give a continuous
map $w$ from the parameter cube to this join, with its literal
affine boundary section. Apply the finite vector coordinate
clocks to $(\mathrm{id},w)$, after a fixed standard scaling of
the target coordinates if needed. The same bounded-variation
calculation gives a standard bound on $w\chi$. Every boundary
coordinate line is affine, so its variation is $ct$ and the
normalized clock $(t+ct)/(1+c)$ is exactly $t$. Hence
$\chi$ fixes the parameter boundary pointwise.

It follows that the map $g$ supplied by A.157 in this
application fixes the boundary and has STANDARD graph
Hausdorff control for $\mathcal G\circ g$. Let $L_\Gamma$
be an ordinary integer bounding that control in the original
untagged graph coordinates, and let $K$ be the common standard
spatial bound of the factors. A point $(x,r_{g(p)}(x))$ can
be matched, by definable compactness, to
$(x',r_{g(q)}(x'))$ at graph distance at most
$L_\Gamma\|p-q\|_\infty$. Therefore
\[
 \|r_{g(p)}(x)-r_{g(q)}(x)\|_\infty
 \le(1+K)L_\Gamma\|p-q\|_\infty.
 \tag{159.10}
\]
The same estimate holds for $s$. No estimate for $r^{-1}$ or
$s^{-1}$ was used. We have obtained a pair of factors with
standard forward parameter bounds and all the equality
conditions of the pair criterion.

Apply that criterion to $r\circ g,s\circ g$. Their quotient
$U\circ g$ is definably homotopic through based concordance
families to a jointly PL family. The convex parameter homotopy
\[
 U_{(1-t)p+t g(p)},\qquad 0\le t\le1,
 \tag{159.11}
\]
is based, fixes the bottom, and connects the original family
to $U\circ g$. This proves (159.1). Translation by PL
representatives handles other base components, and A.113
supplies the component bijection. Empty bases have only their
empty maps; the constructions also apply componentwise and
to zero-dimensional bases.

\textbf{The remaining relative distinction.}\quad
For a nonconstant prescribed PL parameter boundary over $M$,
its parameter Lipschitz constant can be larger than every
ordinary integer. Appending that boundary as an exact section
in A.157 then need not preserve a standard bound. Neither
A.158 nor the present proof erases this exact-boundary issue.
Consequently this theorem does not infer nonreal injectivity
from A.156. A separate parameter-preserving PL field-invariance
argument, or a direct exact-relative standard-bound reduction,
is still needed for that inference. A definable path of PL
coefficient data alone is not such an argument.
\end{proof}



